\chapter{Rencana Selanjutnya}

\section{Rencana Dari Laporan September}

Rencana lanjutan yang tercatat pada laporan progress September adalah:

\begin{enumerate}
    \item Integrasi penuh mapping actor per aksi ke pipeline Kafka Consumer dan Normalizer.
    \item Pengujian end-to-end terhadap klien dengan struktur tabel actor yang berbeda-beda.
    \item Penyempurnaan dokumentasi API dan panduan onboarding klien baru.
    \item Review keamanan pra-produksi, termasuk CORS whitelist, validasi JWT secret, dan konfigurasi akses.
\end{enumerate}

\roadmapdiagram

\section{Rencana Update Terbaru}

Berdasarkan review terbaru, roadmap dapat dibagi menjadi tiga prioritas.

\subsection{Prioritas P0}

\begin{itemize}
    \item Menyelesaikan integrasi actor mapping per operasi ke Kafka consumer dan normalizer.
    \item Menyimpan actor mentah dan actor display secara terpisah.
    \item Membuat endpoint backend khusus \texttt{/api/dashboard/my-users/activity}.
    \item Menjalankan testing Actor Flow V2 menggunakan tabel actor-aware.
    \item Menambahkan pipeline health monitoring dasar, seperti heartbeat, connector status, last event time, dan consumer lag.
\end{itemize}

\subsection{Prioritas P1}

\begin{itemize}
    \item Menambahkan storage observability untuk audit database, Kafka, Docker log, dan WAL database client.
    \item Menyediakan configurable retention policy per client.
    \item Menambahkan dry-run purge untuk evidence yang sudah eligible.
    \item Menambahkan rule-based security alert untuk mass delete, unknown actor, off-hours activity, dan connector down.
    \item Menyusun monitored table reconciliation agar tabel baru dapat ditambahkan tanpa full reinstall.
\end{itemize}

\subsection{Prioritas P2}

\begin{itemize}
    \item Archive evidence ke object storage atau media penyimpanan WORM.
    \item Approval workflow dua pihak untuk perubahan konfigurasi sensitif.
    \item Webhook atau email alert.
    \item Correlation dengan login/auth event, IP, device, dan MFA.
    \item Sequence manifest untuk memperkuat pembuktian completeness per interval.
\end{itemize}

\section{Testing Actor Flow V2}

Testing Actor Flow V2 disiapkan untuk memvalidasi actor resolution dari database client sampai dashboard. Skenario ini menggunakan database yang actor-aware, yaitu database yang memang menyediakan kolom pelaku perubahan.

\begin{longtable}{p{4cm}p{9cm}}
\toprule
\textbf{Skenario} & \textbf{Expected Result} \\
\midrule
INSERT order & Actor terbaca dari field insert actor atau \texttt{created\_by}. \\
\midrule
UPDATE order & Actor terbaca dari field update actor atau \texttt{updated\_by}. \\
\midrule
DELETE atau soft delete & Actor terbaca dari field delete actor atau \texttt{deleted\_by}. \\
\midrule
User source masuk dari CDC & Data user muncul di Web Users atau user tracking. \\
\midrule
Actor berupa UUID & Sistem mencoba resolve UUID ke email atau full name dari \texttt{client\_users}. \\
\midrule
Verify log & Hash dan verification tetap berhasil selama data tidak berubah. \\
\bottomrule
\end{longtable}

\section{Acceptance Criteria}

Core sistem dapat dianggap siap untuk demo lanjutan jika memenuhi kriteria berikut:

\begin{enumerate}
    \item Event INSERT, UPDATE, dan DELETE dari database client muncul di dashboard.
    \item Actor tidak lagi \texttt{Unknown} jika field mapping benar dan data actor tersedia di payload CDC.
    \item Hash verification berhasil untuk log yang tidak berubah.
    \item Merkle proof cocok dengan Merkle root.
    \item Dashboard dapat menampilkan status valid, tampered, pending, dan unreachable.
    \item Admin dapat mendaftarkan client dan mengatur mapping actor melalui dashboard.
\end{enumerate}
